\chapter{Quantum Chromodynamics}

\begin{remarkbox}
Quantum chromodynamics is the non-Abelian gauge theory of the strong
interaction. Its gauge group is \(SU(3)\), its matter fields are quarks, and its
gauge bosons are gluons. QCD is weakly coupled at short distances and strongly
coupled at long distances.
\end{remarkbox}

\section{QCD as an \texorpdfstring{\(SU(3)\)}{SU(3)} Gauge Theory}

QCD is a Yang--Mills theory with gauge group
\[
    SU(3)_{\mathrm{color}}.
\]
The fundamental charges are called color charges. Quarks transform in the
fundamental representation of \(SU(3)\), while gluons transform in the adjoint
representation.

The QCD Lagrangian is
\[
    \mathcal{L}_{\mathrm{QCD}}
    = -\frac{1}{4}G_{\mu\nu}^aG^{a\mu\nu}
      +\sum_f \bar q_f(i\gamma^\mu D_\mu-m_f)q_f.
\]
Here \(f\) labels quark flavors, such as \(u,d,s,c,b,t\). The covariant
derivative is
\[
    D_\mu=\partial_\mu+ig_s G_\mu^a T^a,
\]
where \(g_s\) is the strong coupling, \(G_\mu^a\) are the gluon fields, and
\(T^a\) are \(SU(3)\) generators in the fundamental representation.

\section{Color Charge}

Color is the charge of the strong interaction. A quark field carries a color
index:
\[
    q_i,
    \qquad
    i=1,2,3.
\]
The three values are often called red, green, and blue, but these names are only
labels for basis vectors in an internal vector space.

Gluons carry adjoint color indices:
\[
    G_\mu^a,
    \qquad
    a=1,\ldots,8.
\]
There are eight gluons because the dimension of the adjoint representation of
\(SU(3)\) is eight.

Physical isolated states are color singlets. Examples include mesons,
\[
    \bar q_i q_i,
\]
and baryons,
\[
    \epsilon^{ijk}q_i q_j q_k.
\]
This singlet structure is closely related to confinement.

\section{The Gluon Field Strength}

The QCD field strength is
\[
    G_{\mu\nu}^a
    = \partial_\mu G_\nu^a-\partial_\nu G_\mu^a
      -g_s f^{abc}G_\mu^bG_\nu^c,
\]
up to sign conventions for the covariant derivative. The nonlinear term is the
source of gluon self-interactions.

The pure gauge part of the Lagrangian,
\[
    -\frac{1}{4}G_{\mu\nu}^aG^{a\mu\nu},
\]
contains both three-gluon and four-gluon vertices. This is the main structural
difference between QCD and QED. Photons do not carry electric charge, but gluons
carry color charge.

\section{Gauge Fixing and Ghosts}

Perturbative QCD requires gauge fixing. In a covariant gauge, one adds
\[
    \mathcal{L}_{\mathrm{gf}}
    = -\frac{1}{2\xi}(\partial^\mu G_\mu^a)^2.
\]
The Faddeev--Popov ghost Lagrangian is
\[
    \mathcal{L}_{\mathrm{gh}}
    = -\bar c^a\partial^\mu D_\mu^{ab}c^b,
\]
where
\[
    D_\mu^{ab}=\delta^{ab}\partial_\mu+g_s f^{acb}G_\mu^c.
\]
Unlike in QED, ghosts do not decouple. They interact with gluons and are needed
for consistency of loop calculations.

BRST symmetry organizes the gauge-fixed QCD action and ensures that physical
observables are gauge independent.

\section{Basic QCD Feynman Rules}

The quark propagator is
\[
    \frac{i(\gamma^\mu p_\mu+m_f)}{p^2-m_f^2+i\epsilon}\delta_{ij}.
\]
The gluon propagator in Feynman gauge is
\[
    \frac{-i\eta_{\mu\nu}\delta^{ab}}{k^2+i\epsilon}.
\]
The quark--gluon vertex is
\[
    -ig_s\gamma^\mu T^a_{ij}.
\]
In addition, QCD has three-gluon, four-gluon, and ghost--gluon vertices.

Every QCD diagram has a Lorentz part, a spinor part if quarks are present, and a
color part. The color algebra is as important as the momentum and spinor algebra.
Useful identities include
\[
    \operatorname{tr}(T^aT^b)=\frac{1}{2}\delta^{ab}
\]
and
\[
    [T^a,T^b]=if^{abc}T^c.
\]

\section{Quark--Quark Scattering}

At high energies, quark scattering can be treated perturbatively. A basic process
is quark--quark scattering through one-gluon exchange:
\[
    q q \longrightarrow q q.
\]
The amplitude contains the structure
\[
    [\bar u(p_3)(-ig_s\gamma^\mu T^a)u(p_1)]
    \frac{-i\eta_{\mu\nu}}{q^2+i\epsilon}
    [\bar u(p_4)(-ig_s\gamma^\nu T^a)u(p_2)].
\]
This resembles electron--muon scattering in QED, but with color matrices and the
strong coupling replacing electric charge.

For identical quarks, exchange diagrams and fermionic signs must be included.
For physical hadrons, quarks are confined, so this partonic calculation is only
one ingredient in a full prediction.

\section{Gluon Self-Interactions}

Gluon self-interactions are responsible for much of QCD's distinctive behavior.
The three-gluon vertex is proportional to
\[
    g_s f^{abc},
\]
and the four-gluon vertex is proportional to
\[
    g_s^2 f^{abe}f^{cde}
\]
with permutations of Lorentz and color indices.

Because gluons carry color charge, they interact with other gluons. This creates
nonlinear field dynamics even in the absence of quarks.

The same self-interactions that complicate calculations also drive asymptotic
freedom and confinement. They make QCD qualitatively different from QED.

\section{Asymptotic Freedom}

For QCD with \(N_c=3\) colors and \(n_f\) active quark flavors, the leading beta
function is
\[
    \beta(g_s)
    = -\frac{g_s^3}{16\pi^2}
      \left(11-\frac{2}{3}n_f\right)+\cdots.
\]
For \(n_f<17\), the coefficient is positive inside the parentheses, so the beta
function is negative. Therefore the strong coupling decreases at high energies.

In terms of
\[
    \alpha_s(\mu)=\frac{g_s(\mu)^2}{4\pi},
\]
one finds approximately
\[
    \alpha_s(\mu)\sim \frac{1}{\ln(\mu^2/\Lambda_{\mathrm{QCD}}^2)}.
\]
At large \(\mu\), the coupling is small. This is why perturbative QCD works for
high-energy processes.

\section{Dimensional Transmutation in QCD}

Classically, massless QCD has a dimensionless coupling. Quantum mechanically,
the running coupling introduces a scale
\[
    \Lambda_{\mathrm{QCD}}.
\]
This is dimensional transmutation. Instead of specifying a dimensionless bare
coupling, the low-energy theory is characterized by a physical scale.

The scale \(\Lambda_{\mathrm{QCD}}\) sets the order of hadronic masses and the
onset of strong coupling. It is not the mass of a single particle. It is the
scale at which the perturbative coupling becomes large and nonperturbative
physics becomes important.

\section{Confinement}

At long distances, QCD becomes strongly coupled. Isolated colored particles are
not observed. Instead, physical states are color singlets: mesons, baryons,
glueballs, and more complicated hadronic states.

A qualitative picture is that the potential between heavy quarks grows roughly
linearly at large separation:
\[
    V(r)\sim \sigma r,
\]
where \(\sigma\) is the string tension. Pulling quarks apart stores energy in the
color flux tube until it becomes energetically favorable to produce new quark--
antiquark pairs.

Confinement is not derived by ordinary weak-coupling perturbation theory. It is a
nonperturbative phenomenon, studied with lattice QCD and other methods.

\section{Chiral Symmetry and Its Breaking}

For light quarks, QCD has approximate chiral symmetries. If the quark masses are
neglected, left- and right-handed quarks can transform independently:
\[
    q_L\longrightarrow L q_L,
    \qquad
    q_R\longrightarrow R q_R.
\]
The QCD vacuum does not preserve the full chiral symmetry. Instead, it develops a
quark condensate,
\[
    \langle\bar q q\rangle\neq0,
\]
which breaks chiral symmetry spontaneously.

The associated light pseudoscalar mesons, such as pions, are interpreted as
pseudo-Goldstone bosons. Their small masses reflect the fact that light quark
masses explicitly but weakly break chiral symmetry.

\section{Jets and Partons}

At high energies, asymptotic freedom allows quarks and gluons to behave like
weakly interacting partons during short-distance scattering. However, colored
partons cannot appear as isolated final particles. They hadronize into jets:
collimated sprays of color-singlet hadrons.

A collider event with jets is therefore interpreted in two stages:
\[
    \text{short-distance parton process}
    \longrightarrow
    \text{parton shower}
    \longrightarrow
    \text{hadronization into jets}.
\]
The short-distance process can often be computed perturbatively. The shower is
partly perturbative. Hadronization is nonperturbative and must be modeled or
extracted from data.

\section{Factorization}

Hadron collisions involve bound states, not free quarks and gluons. The
factorization idea separates short-distance partonic scattering from long-
distance hadronic structure.

A schematic factorized cross section is
\[
    \sigma(pp\to X)
    = \sum_{i,j}\int\dd x_1\dd x_2\,
      f_{i/p}(x_1,\mu)f_{j/p}(x_2,\mu)
      \hat\sigma_{ij\to X}(x_1,x_2,\mu).
\]
The functions \(f_{i/p}(x,\mu)\) are parton distribution functions. They encode
the probability of finding a parton \(i\) carrying a fraction \(x\) of the
proton momentum at scale \(\mu\). The short-distance cross section
\(\hat\sigma\) is computed perturbatively.

Factorization is one of the reasons QCD can make predictions for hadron
colliders despite confinement.

\section{Perturbative and Nonperturbative Regimes}

QCD has two faces. At high momentum transfer,
\[
    Q\gg \Lambda_{\mathrm{QCD}},
\]
the coupling is small and perturbation theory is useful. At low momentum
transfer,
\[
    Q\sim \Lambda_{\mathrm{QCD}},
\]
the coupling is large and perturbation theory fails.

Nonperturbative QCD includes confinement, hadron masses, chiral symmetry
breaking, and the detailed structure of hadrons. Perturbative QCD describes
short-distance scattering, high-energy jets, hard emissions, and the scale
evolution of parton distributions.

A realistic QCD prediction often combines both regimes through factorization,
phenomenological inputs, and numerical methods.

\section{Lattice QCD}

Lattice QCD defines the theory nonperturbatively by discretizing spacetime. Gauge
fields are represented by group-valued link variables, and the path integral is
evaluated numerically in Euclidean signature.

The lattice spacing \(a\) acts as a UV regulator. The continuum theory is
recovered by taking
\[
    a\to0
\]
while tuning parameters appropriately. Lattice QCD can compute hadron masses,
decay constants, matrix elements, and other nonperturbative quantities.

Although lattice QCD is technically demanding, it provides the most direct
nonperturbative definition of QCD currently used in practice.

\section{QCD and the Standard Model}

QCD is one of the three gauge sectors of the Standard Model. Its gauge group is
part of
\[
    SU(3)_c\times SU(2)_L\times U(1)_Y.
\]
The \(SU(3)_c\) factor is color. Quarks carry color charge; leptons do not.
Gluons couple only to colored particles.

The strength and structure of QCD shape much of observed matter. Most of the
mass of protons and neutrons comes not from the bare quark masses, but from QCD
binding energy and dynamics.

\begin{summarybox}
QCD is the \(SU(3)\) non-Abelian gauge theory of quarks and gluons. Quarks carry
fundamental color charge, gluons carry adjoint color charge, and physical
isolated states are color singlets. Gluon self-interactions make the theory
asymptotically free at short distances and strongly coupled at long distances.
Confinement, chiral symmetry breaking, jets, factorization, and lattice QCD are
central features of the theory.
\end{summarybox}

\section{Chapter Checklist}

After reading this chapter, one should be able to:
\begin{itemize}
    \item write the QCD Lagrangian;
    \item explain color charge and the role of \(SU(3)\);
    \item distinguish quarks and gluons by representation;
    \item write the gluon field strength;
    \item identify the basic QCD Feynman rules;
    \item explain why gluons self-interact;
    \item state the leading qualitative form of the QCD beta function;
    \item explain asymptotic freedom and dimensional transmutation;
    \item describe confinement and color singlets;
    \item explain the idea of jets and partons;
    \item state the basic meaning of factorization;
    \item distinguish perturbative and nonperturbative QCD regimes.
\end{itemize}
